\section{Follow-on UK data architecture and access}
\label{sec:followon}

This section identifies the data infrastructure needed for a longer UK inheritance programme. \textbf{It is separate from the three Turing DSG deliverables.} Sections \ref{sec:ifs}, \ref{sec:outcomes} and \ref{sec:gifts} turn the data possibilities into distinct empirical research streams.

\subsection{Units of observation}
An ideal research resource would distinguish five records, rather than treating \emph{an inheritance} as one observation:
\begin{enumerate}[nosep]
\item \textbf{Estate/donor:} deceased or living giver, estate value, asset mix, date, taxes, number and relationships of potential recipients.
\item \textbf{Intended transfer:} specific and residuary provisions in a valid will, gifts, trusts, contingent conditions and charitable beneficiaries.
\item \textbf{Realised transfer:} recipient, amount, asset type, net value and actual payment or title-transfer date.
\item \textbf{Recipient background:} age, wealth, education, income, family background, location and prior outcomes.
\item \textbf{Recipient trajectory:} repeated observations of wealth, property, debt, work, business ownership, saving and consumption before and after the transfer.
\end{enumerate}
No single readily accessible UK dataset is known to provide this complete chain. Each source must be evaluated for the subset it genuinely contains.

\subsection{Candidate source families}
\begin{description}[style=nextline]
\item[ELSA (UK Data Service 5050)]
Inheritance histories, parental mortality and repeated economic measures for older people; the main timing and proxy-validation route, subject to variable and licence audits.
\item[Wealth and Assets Survey (WAS; UKDS 7215)]
Wealth stocks, composition, inheritance and inter vivos transfer questions; candidate for inequality and post-receipt balance-sheet analysis. Check exact waves, identifiers, survey weights and transfer purpose before specifying regressions.
\item[Understanding Society (UKDS 6614)]
Possible repeated employment, income, housing and transfer observations, subject to inheritance variables and sufficient recipient events.
\item[ONS Longitudinal Study of England and Wales]
A restricted 1\% population sample linking census and vital events, useful for intergenerational background and cohort structure. The IFS inheritance projections combined this with ELSA and WAS; access is not open.
\item[HMRC Datalab, HMCTS and other administrative custodians]
Possible estate-tax, probate and economic records. Their contents, linkage keys and approval routes are unverified. HMRC's scoping response asks for research questions tied to its functions; access would require a distinct approved project.
\item[Estate administration and private providers]
Potential wills, estate accounts, actual distributions and beneficiary relationships. A partner-specific feasibility study is needed; registration of a will is not proof of access to its text or to actual payments.
\item[Companies House, HM Land Registry and other registers]
Possible business and property outcomes. Public register entries do not by themselves yield accurate person-level links to survey respondents or inheritance recipients.
\end{description}

\subsection{Minimum data dictionary}
For every candidate source, record:
\begin{itemize}[nosep]
\item population and unit of observation; date coverage; geographic coverage and sampling/selection;
\item transfer type (death inheritance, lifetime gift, loan, trust, business succession, charity) and whether it is intended or realised;
\item date, value, gross/net distinction, recipient relationship, siblings/number of heirs and asset composition;
\item longitudinal recipient outcomes and available baseline wealth/background;
\item identifier availability, researcher permissions, confidentiality, release conditions and disclosure controls;
\item uncertainty, missingness, recall error, valuation choices and potential double counting.
\end{itemize}

\subsection{Access sequence and boundaries}
\begin{enumerate}[nosep]
\item Complete an accessible-data audit of WAS and ELSA; start with analyses wholly inside each survey.
\item Seek UK will and provider partnerships to validate distributions, clauses and estate-administration fields in separate, governed samples.
\item Build descriptive external comparisons between samples (for example, shares of unequal bequests), without attempting to re-identify survey participants.
\item Only then scope a custodian-led administrative linkage with specific outcome variables, documented legal basis, approvals and a secure research environment.
\end{enumerate}

Published or licence-bound pseudonymised survey microdata must never be matched to named probate records by reverse-identification. Pseudonymisation is not anonymisation. A future administrative linkage requires approval and custodian-controlled methods. The most valuable data may be \emph{final estate distribution accounts}, rather than will text alone.
